%! Summary Statistics for a Quantitative Variable — Unit 1, Lesson 1.5 — AP Practice
% GRADUAL RELEASE: AP-format practice, assigned but NOT scored — the cover's score
% column reads NA for this component. The graded practice is the `homework` component that
% follows it at the back of the packet.
% Free-response (a) is the spiral to Lesson 1.4 (name the shape from the display) and (b) is
% Skill 4.B in the form the exam asks it — the CHOICE of summary must be justified by the shape.
% MC item 1 is the crux as a multiple-choice item: skewed right with an outlier, so median and
% IQR — and the justification is the reason, not the pair.
% Page lockstep with ap_practice_key rests on \writelines{n} in the blank matching a terse
% \ansline{} in the key — keep key answers short enough not to wrap past the reserved lines.
% The next-lesson preview is NOT here: it closes the packet, in ../homework/.
%
% THE DATA (15 households, gallons per day):
%   110 120 130 140 150 160 180 | 200 | 220 240 270 310 350 400 620
%   sum 3600 -> mean 240 ; median 200 ; Q1 140 ; Q3 310 ; IQR 170 ; range 510 ; s_x ~ 136.6
%   Right-skewed with one large value at 620. Both outlier rules flag the 620:
%   Q3 + 1.5(IQR) = 310 + 255 = 565 < 620, and mean + 2s = 240 + 273.2 = 513.2 < 620.
\documentclass[12pt]{article}
\usepackage{apstats-article}
\usepackage{apstats-boxes}

\begin{document}

\pageheader{Unit 1, Lesson 1.5}{AP Practice}
% NAMESTRIP: no \namedateperiod here — the name row lives on the cover only.

\vspace{0.10in}

% Authored BYTE-IDENTICALLY in ap_practice/ and ap_practice_key/ — this box is what keeps the
% blank and the key the same number of pages.
\boxguard[8]
\begin{remindbox}
\small
This page is \textbf{not required and not scored} --- your graded homework is the last section
of this packet. Work these AP-format reps and bring what stalls you to study hall.
\end{remindbox}

\vspace{0.08in}

\boxguard[30]
\begin{scenariobox}[Household Water Use]{navy}
\small
A water utility recorded each of 15 homes' average daily water use, in gallons. The stem is
the \emph{hundreds} digit and each leaf the \emph{tens} digit: row \textbf{3} is 310 and 350.

\vspace{3pt}
\begin{center}
\renewcommand{\arraystretch}{1.2}
\begin{tabular}{@{} r | l @{}}
\textbf{1} & 1 \quad 2 \quad 3 \quad 4 \quad 5 \quad 6 \quad 8 \\
\textbf{2} & 0 \quad 2 \quad 4 \quad 7 \\
\textbf{3} & 1 \quad 5 \\
\textbf{4} & 0 \\
\textbf{5} & \\
\textbf{6} & 2 \\
\end{tabular}
\end{center}
\vspace{3pt}

Computer output for the same 15 values:

\vspace{2pt}
\renewcommand{\arraystretch}{1.2}
\begin{tabularx}{\linewidth}{@{} l Y Y Y Y Y Y Y @{}}
\rowcolor{sky}
\textbf{$n$} & \textbf{Mean} & \textbf{SD} & \textbf{Min} & \textbf{$Q_1$} & \textbf{Median} &
\textbf{$Q_3$} & \textbf{Max} \\
\hline
15 & 240.0 & 136.6 & 110 & 140 & 200 & 310 & 620 \\
\end{tabularx}
\end{scenariobox}

\vspace{0.10in}

\boxguard[16]
\begin{headlinebox}{goldbg}
\small
\textbf{Section I --- Multiple Choice.} Circle the single best answer. Every answer choice is
about the context, not just the arithmetic --- read them all before choosing.
\end{headlinebox}

\vspace{0.10in}

\begin{enumerate}[leftmargin=1.9em, itemsep=0.10in]

  \item Which pair of summary statistics should be reported for these 15 homes, and why?
        \begin{enumerate}[label=\textbf{(\Alph*)}, leftmargin=2.2em, itemsep=2pt]
          \item The mean and standard deviation, because they use all 15 values.
          \item The mean and the IQR, because the mean is the more accurate center.
          \item The median and the IQR, because the distribution is skewed to the right with one
                home far above the rest, and both of those summaries are resistant.
          \item The median and the standard deviation, because the median is resistant.
          \item The range and the IQR, because both of them measure variability.
        \end{enumerate}

  \item The 620-gallon home is re-metered and its correct figure is 1{,}240 gallons per day
        --- double the recorded value. Which summary statistic changes \textbf{by the most}?
        \begin{enumerate}[label=\textbf{(\Alph*)}, leftmargin=2.2em, itemsep=2pt]
          \item The median
          \item The interquartile range
          \item The mean
          \item The range
          \item None of them change, because only one home was corrected.
        \end{enumerate}

  \item The utility converts all 15 values to litres by multiplying each by 3.8. Which is
        true of the summary statistics for the new list?
        \begin{enumerate}[label=\textbf{(\Alph*)}, leftmargin=2.2em, itemsep=2pt]
          \item Only the mean is multiplied by 3.8; the others are unchanged.
          \item The mean and median are multiplied by 3.8, but the standard deviation is unchanged.
          \item The mean is multiplied by 3.8 and the standard deviation by $3.8^2$.
          \item None of them change, because a change of units does not affect a statistic.
          \item The mean, median, range, IQR, and standard deviation are all multiplied by 3.8.
        \end{enumerate}

  \item A city-wide report states that a household using 350 gallons per day is at the
        \textbf{90th percentile} of daily water use for the city. This means that
        \begin{enumerate}[label=\textbf{(\Alph*)}, leftmargin=2.2em, itemsep=2pt]
          \item 90\% of the city's households use more than 350 gallons per day.
          \item the household uses 90\% of the water in its neighbourhood.
          \item the household uses 90 gallons per day more than the city average.
          \item about 90\% of the city's households use 350 gallons per day or less.
          \item 350 gallons is 90\% of the largest amount used by any household in the city.
        \end{enumerate}

\end{enumerate}

\vspace{0.10in}

\boxguard[16]
\begin{headlinebox}{goldbg}
\small
\textbf{Section II --- Free Response.} Show your reasoning. Every answer must be written
\emph{in the context of the problem} --- that is what earns credit on the AP exam.
\end{headlinebox}

\vspace{0.10in}

\begin{enumerate}[leftmargin=1.9em, itemsep=0.12in, resume]

  \item Use the water-use record above. \textbf{(a)} Describe the shape of this distribution of
        daily water use, and identify any value that stands apart from the rest.\\[4pt]
        \writelines{2}

        \vspace{4pt}
        \textbf{(b)} Report \emph{one} measure of center and \emph{one} measure of variability
        for these 15 homes. Justify your choice of pair using the shape of the distribution.\\[4pt]
        \writelines{2}

        \vspace{4pt}
        \textbf{(c)} Interpret $Q_1 = 140$ gallons per day in the context of this street.\\[4pt]
        \writelines{2}

        \vspace{4pt}
        \textbf{(d)} A resident tells a neighbour, ``A typical home on this street uses 240
        gallons a day.'' Explain why that overstates typical use here, naming a specific feature
        of the distribution.\\[4pt]
        \writelines{2}

\end{enumerate}



\end{document}
